% 来源：OEBPS/Text/01_44.xhtml（OPF spine 顺序；仅第一册）
\cleardoublepage
\BookEntry[富兰克林《富兰克林自传》]{富兰克林《富兰克林自传》}

\BookCover{images/书单28富兰克林自传.jpg}{[美] 富兰克林著| 蒲隆译| 译林出版社2015年版}

\begin{BookSummaryBox}
\textbf{美国迄今为止最重要、读者最多的自传作品之一}
\par
\end{BookSummaryBox}

\subsubsection*{这本书谈什么？}

富兰克林是一名政治家，也是杰出的科学家、外交家、出版家、作家和社会事业家，他像是“从天上偷窃火种的第二个普罗米修斯”，是美国的国父、美国独立运动领袖、民主精神的缔造者，及《独立宣言》的起草人。他已成为举世公认的现代文明之父、美国人的象征。
\par

在《富兰克林自传》，他写下了关于人生奋斗与成功的真知灼见，以及诸种善与美的道德真谛，被公认为是改变了无数人命运的美国精神读本。整部自传在通俗易懂的叙述中不时会有睿智和哲理的火花，开创了美国传记文学的优良传统，使自传成为一种新的文学体裁。世界各地的青年人都希望学习富兰克林成功的秘诀，他们把这部书当成“人生指导”读物。
\par

\subsubsection*{芒格为什么谈起这本书？}

富兰克林对芒格的影响实在太大了。这本《穷查理宝典》效仿的对象便是本杰明·富兰克林的《穷理查年鉴》。查理喜欢思考人类思维的错误，他在人类误判心理学的演讲之前便说道：
\par

\begin{quotation}
对我帮助极大的本杰明·富兰克林给世人留下了他的自传和其他许多东西。
\par
\end{quotation}

芒格在《不那么沉默的合伙人》一文中，还提到：
\par

\begin{quotation}
我曾经蹩脚地试图模仿本杰明·富兰克林的生活。富兰克林在42岁那年放弃经商，专注于成为作家、政治家、慈善家、投资者和科学家。所以我也将我的兴趣从生意场上移开。
\par
\end{quotation}

\subsubsection*{《富兰克林自传》作者富兰克林如是说}

“我未曾见过一个早起、勤奋、谨慎、诚实的人抱怨命运不好。良好的品格，优良的习惯，坚强的意志，是不会被假设所谓的命运打败的。”
\par

“对所有的人以诚相待，同多数人和睦相处，和少数人常来常往，只跟一个人亲密无间。”
\par

“一个人失败的最大原因，就是对于自己的能力永远不敢充分信任，甚至自己认为必将失败无疑。”
\par

“人生的悲剧就在于我们衰老得太早而又聪明得太晚。”
\par

“把别人对你的诋毁放在尘土中，而把别人对你的恩惠刻在大理石上。”
\par

“惟坚韧者始能遂其志。”
\par
